\section*{Chapter \ref{ch:m2:getting-access-rates-credit} --- Getting Access: Rates and Credit}

\begin{solution}{exo:m2:getting-access-rates-credit:1}
On a default it replaces every outstanding trade between the two parties by a single
amount owed one way, so the survivor loses only the net value. If the netting is not
enforceable where the client is, a liquidator could claim the trades the dealer owes
while leaving it to claim, as an unsecured creditor, those the client owes: the
dealer's exposure, and the capital held against it, would be gross.
\end{solution}

\begin{solution}{exo:m2:getting-access-rates-credit:2}
An LEI; an \omterm{def:m2:getting-access-rates-credit:isda}{ISDA master agreement}, schedule and \omterm{def:m2:swap-clearing-and-the-ccp-basis:umr}{credit support annex} with each executing
dealer; account and clearing agreements with a swap \omterm{def:m2:getting-access-rates-credit:cb}{clearing broker}; onboarding with the
execution platforms; and confirmations for each trade.
\end{solution}

\begin{solution}{exo:m2:getting-access-rates-credit:3}
A single legal entity, by a 20-character code; its reference data answer who is who and
who owns whom.
\end{solution}

\begin{solution}{exo:m2:getting-access-rates-credit:4}
$B^* = (8 - 0.5)\text{ million} / (11.5 - 5.9)\text{ bp} = 7.5\text{ million} / 0.00056 =
\text{USD}~13.4$ billion.
\end{solution}

\begin{solution}{exo:m2:getting-access-rates-credit:5}
Sponsored clearing ties the fund to its sponsor, which trades with it or submits its
trades and guarantees them; a done-away agent clears trades the fund does with anyone,
so the fund can choose its trading counterparties apart from its clearing member. The
agent still guarantees the fund to the clearing house and charges for it.
\end{solution}

\begin{solution}{exo:m2:getting-access-rates-credit:6}
Once eligible repo must be cleared, from June 2027, much of the bilateral route is no
longer available whatever it costs; the choice is between the cleared routes, and the
bilateral costs matter only for trades outside the mandate.
\end{solution}

\begin{solution}{exo:m2:getting-access-rates-credit:7}
Without swaps the break-even is USD~15.0 billion: the swaps' lower fees at the direct
route no longer help pay the fixed cost. With membership at 12 million it is USD~20.5
billion: the extra fixed cost must be paid from the same saving per dollar.
\end{solution}

\begin{solution}{exo:m2:getting-access-rates-credit:8}
The fixed costs recur every year: staff, systems, legal and the capital in the clearing
fund. Direct membership is cheaper only above the break-even size, which depends on the
spread sponsors charge; it also brings obligations the model leaves out, such as
loss-sharing and default management, and it cannot be switched off quickly if the book
shrinks.
\end{solution}

\begin{solution}{pb:m2:getting-access-rates-credit:1}
\textbf{1.} Rent 11.5, buy 5.9, bilateral 18.5 basis points a year per dollar of repo.
\textbf{2.} USD~6.25 million, 10.95 million and 9.45 million a year.
\textbf{3.} Renting.
\textbf{4.} Below USD~0.43 billion of repo.
\textbf{5.} It is capital the fund cannot use elsewhere, and it is at risk if another
member defaults.
\textbf{6.} USD~13.4 billion.
\textbf{7.} USD~15.9 million a year.
\textbf{8.} USD~3.7 million a year.
\textbf{9.} It rises to USD~15.0 billion.
\textbf{10.} It rises to USD~20.5 billion.
\textbf{11.} USD~28.8 billion at 5 basis points; USD~7.8 billion at 12.
\textbf{12.} It depends on the sponsor's balance sheet, its capital rules and appetite,
and the fund's size and behaviour, and it can change in a stress.
\textbf{13.} Default-fund loss sharing, liquidity calls from the clearing house, default
management duties, and the cost of meeting membership standards as they change.
\textbf{14.} Dependence on the sponsor, which can raise its price or cut the fund's
limits in a stress, and less netting across counterparties than a member gets.
\textbf{15.} Eligible repo must be cleared from June 2027, so the bilateral route
shrinks to what the mandate does not cover.
\textbf{16.} Building takes time (applications, systems, staff), so it should start
before the book passes the break-even, while renting continues, if growth is likely
enough.
\textbf{17.} For flexibility, to avoid fixed commitments and default-fund exposure, or
because the sponsor's netting and balance sheet serve it better than its own.
\textbf{18.} The LEI, master agreements with dealers and repo counterparties, \omterm{def:m2:swap-clearing-and-the-ccp-basis:umr}{credit
support annexes}, and clearing agreements.
\textbf{19.} Named result: \emph{build or rent}: direct membership beats a \omterm{def:m2:getting-access-rates-credit:cb}{clearing
broker} above USD~13.4 billion of repo, with swaps twice as large, where each route costs
USD~15.9 million a year.
\textbf{20.} Lower costs per unit and control of its access, paid for with fixed costs
and a share of the clearing house's risks.
\end{solution}

\begin{solution}{iq:m2:getting-access-rates-credit:1}
The standard contract governing all OTC derivatives between two parties, with a
negotiated schedule and per-trade confirmations forming one agreement. Close-out
netting replaces all trades by one amount on a default, so exposure and capital are
measured net, provided it is enforceable in the counterparty's insolvency.
\iqlookfor{single agreement, netting, enforceability.}
\end{solution}

\begin{solution}{iq:m2:getting-access-rates-credit:2}
As a \omterm{def:m2:getting-access-rates-credit:sponsored}{sponsored member} through a bank sponsoring member of FICC, which submits and
guarantees its trades; through a member's agent clearing (done-away) service; or, for a
very large firm, by direct membership. Bilateral uncleared repo is an alternative that
the \omterm{def:m2:swap-clearing-and-the-ccp-basis:mandate}{clearing mandate} narrows.
\iqlookfor{the routes and the mandate.}
\end{solution}

\begin{solution}{iq:m2:getting-access-rates-credit:3}
List each route's fixed costs and its variable costs per unit of business: fees,
funded margin, capital in clearing funds; compute annual costs as the book grows; the
break-even is the fixed-cost difference over the unit-cost difference; test the
sensitivity to the sponsor's spread and to growth, and add what the model leaves out.
\iqlookfor{fixed versus variable, break-even, sensitivity.}
\end{solution}

\begin{solution}{iq:m2:getting-access-rates-credit:4}
It clears the fund's swaps at the clearing house, posts margin and guarantees the fund,
charging fees and setting limits. In a stress it can raise margin add-ons, cut limits,
ask the fund to reduce positions or move them, and, if the fund defaults, close them out.
\iqlookfor{clearing, margin, limits and their use in stress.}
\end{solution}

\begin{solution}{iq:m2:getting-access-rates-credit:5}
Principal trading firms trade on the inter-dealer order books through the platforms'
onboarding and with a bank or clearing member standing behind their credit and
settlement; in 2014 they were more than half of the inter-dealer volume studied on 15
October.
\iqlookfor{platform access plus credit intermediation.}
\end{solution}

\begin{solution}{iq:m2:getting-access-rates-credit:6}
One record per counterparty and legal entity: documents signed (master agreements,
schedules, annexes, clearing and platform agreements) with versions and dates, the key
terms extracted (eligible collateral, thresholds, termination events, allowed
products), and the permissions they give; expose them to the order path as limits and
to risk as netting sets; alert on expiries, missing documents and changes.
\iqlookfor{structured terms feeding trading permissions.}
\end{solution}
